\section{Stage 3B: observable guarantees and selective-label identifiability}
\label{sec:stage3b-identifiability}

This section freezes every Stage 3A statement and asks a different question:
what follows from the selector and the labels that S6-AE actually records,
without a full-pool prediction-error bound, a solver-drift bound, or an
exchangeability assumption?  The answer has two parts.  The selector admits
an exact, deterministic exchange identity.  In contrast, key true-hardness
quantities are not point identified from selectively acquired labels.

\subsection{Objects, total order, and observability classes}

Fix a cycle $k$ and suppress $k$ where no ambiguity results.  The candidate
union $C$ is finite and hash-deduplicated.  Define the strict total order
$\prec$ by descending HGB40 prediction $\widehat f_k$ and, for equal
predictions, ascending instance hash.  Let $\operatorname{Top}_r(A)$ denote
the first $r$ elements of $A$ in this order.  For S6-AE,
\[
  m=10{,}000,\qquad K=1{,}000,\qquad
  G=\operatorname{Top}_m(C),
\]
where $G$ is the frozen, diagnostic pure-global HGB selection.  The per-niche
floor $Q$ contains one $\prec$-best member of every occupied niche, so
$q=|Q|\le K\le m$.  The actual acquisition and next-parent population is
\[
  T=P_{k+1}
   =Q\cup \operatorname{Top}_{m-q}(C\setminus Q).
\]
All set statements below concern instance identities.  The stored row order of
$T$---floor first and global fill second---is not part of those statements.

True current-checkpoint hardness remains
\[
  f_k(x)=\log\!\bigl(1+L_{BQ}(M_k,x)\bigr).
\]
Let $V_k\subseteq C_k$ be the candidates for which a verified value of $f_k$
under the entering checkpoint $M_k$ is available.  In the inspected S6-AE
artifacts this includes every member of $T_k$ and any current-cycle calibration
member whose hash also lies in $C_k$.  A label obtained under $M_t$, $t\ne k$,
is not thereby a member of $V_k$.

It is useful to separate four information classes.
\begin{enumerate}
  \item \emph{Deterministically observable:} candidate identities, hashes,
  descriptors, niches, HGB predictions and ranks, and the sets $Q,T,G$.
  \item \emph{Verified true-hardness:} $f_k(x)$ only for $x\in V_k$, with
  checkpoint identity and hash alignment verified.
  \item \emph{Latent:} any functional that depends on an unobserved
  $f_k(x)$ for $x\in C_k\setminus V_k$, such as a full-pool niche maximum.
  \item \emph{Hypothetical:} quantities requiring a new sampling law, error
  model, hardness-range restriction, future mutation law, or training-value
  relation.
\end{enumerate}
Write $\mathcal O_k$ for the first two classes.  A finite-pool quantity is
\emph{point identified} from $\mathcal O_k$ if every completion of $f_k$ on
$C_k\setminus V_k$ that agrees with $\mathcal O_k$ gives the same value.

\subsection{Exact Elite-QD versus global-top-$m$ exchange}

\mstatus{PROVED}\par
\istatus{VERIFIED for the inspected S6-AE selector and artifacts}\par
\nstatus{STANDARD CONSTRAINED TOP-$m$ EXCHANGE}\par

\begin{proposition}[Exact protected-floor exchange]
\label{prop:qd-global-exchange}
Let $C$ have a strict total order, let $G=\operatorname{Top}_m(C)$, and let
$Q\subseteq C$ satisfy $q=|Q|\le m$.  If
\[
  T=Q\cup\operatorname{Top}_{m-q}(C\setminus Q),
\]
then
\[
  T\setminus G=Q\setminus G
\]
and
\[
 |T\setminus G|=|G\setminus T|\le |Q|\le K.
\]
\end{proposition}

\begin{proof}
There are $m-|G\cap Q|\ge m-q$ elements of $G\setminus Q$.  Every one of
them precedes every element of $C\setminus G$.  Hence the first $m-q$
elements of $C\setminus Q$ all belong to $G$, so the global-fill part of $T$
contributes nothing to $T\setminus G$.  Therefore
$T\setminus G=Q\setminus G$.  Both $T$ and $G$ have cardinality $m$, so their
two directed set differences have equal cardinality.  The final inequalities
are immediate.
\end{proof}

Thus the actual population differs from pure global selection only through
the floor elites that lie outside the pure top $m$, together with the same
number of displaced global members.  The result does not state that all of
$Q$ is ``forced'': $Q\cap G$ costs no global slot.

Define the predicted-score opportunity cost
\[
  \Delta_k^{\mathrm{QD}}
   =\sum_{x\in G_k}\widehat f_k(x)
    -\sum_{x\in T_k}\widehat f_k(x).
\]

\begin{proposition}[Exact surrogate-score cost]
\label{prop:qd-predicted-cost}
Under Proposition~\ref{prop:qd-global-exchange},
\[
  \Delta_k^{\mathrm{QD}}
  =\sum_{x\in G_k\setminus T_k}\widehat f_k(x)
   -\sum_{x\in Q_k\setminus G_k}\widehat f_k(x)
  \ge 0.
\]
It is exactly computable from the frozen HGB scores and membership records.
\end{proposition}

\begin{proof}
The common members cancel.  The two remaining sets have the same size, and
every member of $G_k$ has score at least that of every candidate outside
$G_k$.  Equivalently, $G_k$ maximizes the sum of scores among all
$m$-subsets of $C_k$.  Hash tie-breaking makes membership unique but permits
$\Delta_k^{\mathrm{QD}}=0$ when boundary scores tie.
\end{proof}

This is a deterministic surrogate-score tradeoff, not true-LBQ regret.  It
does not imply a change in future reachability or downstream training value.
It is an elementary exchange property of a protected set plus top-$m$ fill,
not a new QD theorem.

\subsection{Read-only S6-AE artifact audit}

The selection identities were replayed from the ten frozen cycles without new
solver calls.  Let $r_k=|T_k\setminus G_k|$ and let $D_k=G_k\setminus T_k$.
The label set used in the table is
\[
 V_k=T_k\cup(B_k\cap C_k),
\]
where $B_k$ is the current-checkpoint calibration table.  ``Unlab. $D$'' is
$|D_k\setminus V_k|$.

\begin{center}
\small
\begin{tabular}{crrrrrr}
\toprule
cycle & $|C_k|$ & $|Q_k|$ & $r_k$ & $|V_k|$ & unlab. $D_k$ & $\Delta_k^{\rm QD}$\\
\midrule
C00 & 119781 & 435  & 41  & 10409 & 41  & 0.574\\
C01 & 119700 & 831  & 463 & 18530 & 421 & 4.415\\
C02 & 119507 & 991  & 711 & 14306 & 662 & 9.928\\
C03 & 119487 & 1000 & 800 & 12899 & 765 & 14.322\\
C04 & 119531 & 1000 & 838 & 12231 & 803 & 22.946\\
C05 & 119547 & 1000 & 848 & 11767 & 820 & 27.599\\
C06 & 119667 & 1000 & 853 & 11521 & 821 & 41.746\\
C07 & 119661 & 1000 & 873 & 11358 & 849 & 75.619\\
C08 & 119700 & 1000 & 873 & 11143 & 860 & 70.111\\
C09 & 119664 & 1000 & 876 & 10977 & 862 & 83.503\\
\bottomrule
\end{tabular}
\end{center}

The cost is the sum on the HGB log-target scale; dividing by $m$ gives the
stored mean-score difference.  All ten cycles satisfy the exact set identity.
Once coverage reaches 1,000 niches at C03, $|Q_k|$ remains fixed but $r_k$ and
$\Delta_k^{\mathrm{QD}}$ continue to change.  Coverage saturation therefore
fixes the number of protected constraints, not their overlap with global top
$m$ or their surrogate opportunity cost.

The artifacts suffice for the deterministic columns and for selected-set
truth summaries.  They do not suffice for a true-hardness comparison of $T_k$
and $G_k$: every cycle has unverified members of $D_k$.  In particular, the
stored \texttt{training\_exact\_hardness\_cost.csv} gives its $G$ row only on
naturally labeled members (9,959 at C00 and 9,124 at C09), whereas its $T$ row
has all 10,000 members.  That descriptive difference is not a full-set
$G$-versus-$T$ estimand.

\subsection{Selective-label indistinguishability}

\mstatus{PROVED IMPOSSIBILITY}\par
\istatus{VERIFIED INFORMATION OBSTRUCTION}\par
\nstatus{STANDARD SELECTIVE-LABEL RESULT}\par

\begin{proposition}[Finite-pool non-identifiability]
\label{prop:selective-label-impossibility}
Suppose $C\setminus V$ is nonempty.  Without a restriction relating
$\widehat f$ to $f$, or otherwise restricting unobserved values of $f$, there
exist two finite worlds with identical candidate identities, descriptors,
niches, HGB predictions, $Q,T,G$, and labels on every member of $V$, but
different true-hardness values on $C\setminus V$.

If an unverified candidate $z$ shares niche $j$ with its selected elite, the
two worlds can have arbitrarily different niche maxima and selected-elite
regret.  With sufficiently many unverified nonselected candidates they can
also give top-tail recall one and zero for the same selected set.  If
$(G\setminus T)\setminus V$ is nonempty, the sign and magnitude of the true
hardness sum difference between $G$ and $T$ are not identified.
\end{proposition}

\begin{proof}
Keep the finite pool and every observable fixed.  For the niche claim, set the
verified elite value to one, set every other candidate in that niche to at
most one in world $W_0$, and set $f(z)=H>1$ in world $W_H$.  The selector is
unchanged because it uses $\widehat f$, but the niche regret changes from zero
to $H-1$.  Since $H$ is unrestricted, the difference is unbounded.

For top-$r$ recall, take $r$ selected candidates with fixed verified value
one and $r$ unverified nonselected candidates.  Give the latter value zero in
$W_0$ and value $H>1$ in $W_H$.  With deterministic truth tie-breaking, the
top-$r$ set is entirely selected in $W_0$ and entirely unselected in $W_H$.

Finally let $d\in(G\setminus T)\setminus V$.  Common members of $G$ and $T$
cancel in their true-score sums, while every member of $T\setminus G$ is
verified because $T\subseteq V$.  Setting $f(d)=0$ or $f(d)=H$ leaves all
observables fixed and can reverse the comparison; increasing $H$ makes the
difference arbitrarily large.  Multiple unverified displaced candidates only
strengthen the construction.
\end{proof}

This is an identifiability statement, not evidence that the realized S6-AE
selection is poor.  It says that the inspected observations are compatible
with sharply different unseen true-hardness landscapes.

\subsection{Partial identification under a known hardness range}

Suppose, additionally, that a scientifically justified finite range
$L\le f_k(x)\le U$ is known for every member of the finite candidate pool.
Let
\[
 A=T\setminus G,\qquad D=G\setminus T,
 \qquad D_V=D\cap V,\qquad D_U=D\setminus V.
\]
Because $A\subseteq T\subseteq V$, the following interval is observable:
\begin{align*}
 &\sum_{x\in D_V} f_k(x)+|D_U|L-\sum_{x\in A}f_k(x)\\
 &\qquad\le
 \sum_{x\in G}f_k(x)-\sum_{x\in T}f_k(x)\\
 &\qquad\le
 \sum_{x\in D_V} f_k(x)+|D_U|U-\sum_{x\in A}f_k(x).
\end{align*}
These bounds are sharp absent further restrictions: assign every unobserved
displaced value to $L$ or to $U$.  A known range therefore yields partial,
not point, identification.  Nonnegativity of LBQ supplies $L=0$ on the log
scale, but no inspected artifact supplies a finite scientific $U$.  The
resulting upper endpoint is therefore unavailable, and the read-only lower
endpoints do not determine the sign in any inspected cycle.

For a threshold $\tau$, the acquired hard count
$|\{x\in T:f_k(x)\ge\tau\}|$ is exact.  Full-pool recall is only bounded by
allowing every latent candidate to fall below or above $\tau$.  Exact
full-pool maxima similarly require a census of the relevant niche or an
additional tail/range model.

\subsection{What the existing labels do and do not identify}

The following selected-population quantities are exactly computable without
new true-LBQ calls: the current $f_k$ distribution on $T_k$; the acquired
count above a predeclared threshold; current hardness of each measured
per-niche elite; HGB residuals on $V_k$ or on $T_k$ with the conditioning set
stated; $T_k$--$G_k$ overlap; $r_k$; and $\Delta_k^{\mathrm{QD}}$.

They characterize selected/acquired candidates only.  They do not identify a
full-pool HGB error distribution, a full-pool hard-tail recall, a niche
maximum over $C_k$, historical improvement under the current checkpoint,
generalization to another checkpoint, future reachability, or downstream BQ
training value.  The archive calibration sample is randomized within the
archive, not within the current candidate union; after intersecting it with
$C_k$ it is not a documented probability sample of $C_k\setminus T_k$.
Selected-set residuals therefore cannot be extrapolated to the union without
an explicit selection/sampling model.

\subsection{Minimal additional information}

Different targets require different additions.
\begin{enumerate}
  \item \emph{Exact current $T$ versus $G$ comparison.}  Every member of
  $A=T\setminus G$ is already labeled.  It is enough to evaluate the missing
  members of $D=G\setminus T$ under the same entering checkpoint.  By
  Proposition~\ref{prop:qd-global-exchange}, this requires at most $K=1,000$
  extra labels, and often fewer after reusing verified current labels.  This
  is an exact paired certificate, but it has not been completed in the frozen
  artifacts.
  \item \emph{Finite-pool means, tail mass, or recall.}  Draw unselected
  candidates with known, positive inclusion probabilities, preferably
  stratified by niche and frozen HGB-score bins.  Design-based weighting can
  then estimate declared aggregate quantities.  Deterministic selection
  probabilities of zero cannot support inverse-propensity extrapolation.
  \item \emph{Niche maxima or elite regret.}  Verify the entire target niche,
  verify a declared finite contender set with a valid local error certificate,
  or impose and validate a tail/range model.  A small random sample alone does
  not identify an exact maximum.
  \item \emph{Cross-cycle claims.}  Re-evaluate paired identities under the
  same current checkpoint or posit a separately validated drift model.  Old
  labels are not substitutes.
  \item \emph{Training value.}  A causal training comparison is required;
  acquisition hardness alone supplies no identification.  Stage 3B does not
  propose or run such a comparison.
\end{enumerate}

A minimal future verification protocol, if separately authorized, would
freeze $C,Q,T,G$ before truth evaluation; predeclare the checkpoint, estimand,
strata, and inclusion probabilities; label all still-unlabeled $G\setminus T$
for the exact paired selector comparison; add a probability sample from
$C\setminus(T\cup G)$ only for full-pool aggregate estimands; preserve hash
joins and nonresponse accounting; and report design-based uncertainty.
This is a protocol specification, not an experiment performed here.

\subsection{Novelty and Stage 3B decision}

The protected-floor exchange and its $K$-label paired certificate are useful
implementation audits, but the proof is elementary constrained top-$m$
selection.  QD already formalizes retaining high-quality representatives
across behavioral regions.  Selective-label non-identifiability, partial
identification under bounded outcomes, and propensity-based learning under
partial feedback are established themes
\cite{Lakkaraju2017,Manski2005,SwaminathanJoachims2015}.  Surrogate-assisted
evolutionary computation likewise treats exact-evaluation management as a
core issue rather than equating surrogate scores with truth \cite{Jin2011}.

The BLADE-specific contribution here is a precise audit map: because the
actual $T/G$ symmetric difference is at most $K$, an exact same-checkpoint
comparison needs labels only for the missing displaced set, not the roughly
120,000-member union.  This is practically useful but not, by itself, a
publication-level theorem.  Existing frozen labels do not complete that
certificate and cannot identify the broader full-pool claims.

The Stage 3B research decision is
\[
 \boxed{\text{B -- IDENTIFIABILITY\_LIMIT\_CONFIRMED}.}
\]
No Stage 3C, new true-LBQ call, mutation, evolution, HGB training, BQ update,
or algorithm modification is initiated.
